% ===================================================================
%  TAVOLA N. 4 — Età matura e vecchiaia.
%  Traduction 2026 d'apres texto/io, controlee sur texto/fr et
%  texto/en. Consigne : voir texto/it/10-tavola-01.tex.
%
%  DEUX MOTS DE CETTE PAGE SE REPONDENT, ET C'EST TOUTE UNE SOCIETE
%  QUI TIENT DANS LEUR DIFFERENCE.
%
%  L'alinea 9 dit que ce mariage unit deux familles riches. Le mariage
%  s'appelle en italien matrimonio, et « matrimonio » est le
%  latin « matri-monium », LA CHARGE DE LA MERE. Le mot qui lui
%  ressemble, « patrimonio », est « patri-monium », LA CHARGE DU
%  PERE — et c'est l'heritage, la fortune, ce qui passe de main en
%  main a l'alinea 9 justement.
%
%  LE MARIAGE EST NOMME PAR CE QUE LA FEMME Y PREND, L'HERITAGE PAR CE
%  QUE L'HOMME Y LAISSE. Les deux mots sont batis sur le meme moule, a
%  la meme epoque, par les memes juristes, et ils partagent la page :
%  le premier a l'alinea 9, le second dans tout ce que l'alinea 9
%  raconte. Aucune autre langue des trente-quatre ne porte cette paire.
%
%  ET LES SPOSI DE L'ALINEA 10 SONT « LES PROMIS ». « Sposo » et
%  « sposa » viennent de « sponsus » et « sponsa », participes de
%  « spondere », promettre solennellement — le meme verbe qui a donne
%  « rispondere » et « responsabile ». La jeune femme de l'alinea 10
%  s'appelle donc, mot pour mot, celle a qui l'on a repondu.
%
%  LA SECONDE SCENE REPOND AU TABLEAU 4 ALLEMAND, ET PAR L'ABSENCE.
%
%  L'allemand a deux mots pour chaque maladie — Lungenentzündung
%  contre Pneumonie, Erkältung contre Katarrh, Ohnmacht contre
%  Synkope — et choisir lequel etait, la-bas, un choix de PERSONNAGE :
%  la grand-mere ne dit pas ce que dit son medecin.
%
%  L'italien n'a pas ce choix a faire, et c'est pour la raison
%  inverse. Sa couche savante n'est pas etrangere : « polmonite »,
%  « reumatismi », « vertigini », « febbre », « sincope » sont a la
%  fois les mots du medecin et ceux de tout le monde, parce que le
%  latin de l'un est la langue de l'autre. LA OU L'ALLEMAND DOUBLE,
%  L'ITALIEN FOND — et les deux fois c'est la meme cause : une langue
%  double son vocabulaire savant quand elle l'emprunte, et le confond
%  avec le sien quand elle en est l'heritiere.
%
%  Une seule paire subsiste, et elle est populaire des deux cotes :
%  « raffreddore » et « colpo di freddo ». La colonne ecrit le
%  premier.
% ===================================================================

%%K t04-apar-1 apar
\VUsaut{4.0mm}
\VUtitre{15.0pt}{60}{TAVOLA N. 4}
\VUsaut{1.6mm}
\VUtitre{11.4pt}{0}{Età matura e vecchiaia.}
\VUsaut{3.2mm}

%%K t04-tit-1 sub
\VUcentre{11.4pt}{0}{\textit{Prima scena.}}
\VUsaut{1.6mm}
\VUtitre{11.4pt}{0}{Il pranzo di nozze.}
\VUsaut{2.6mm}

%%K t04-tit-2 sub
\VUpk{10.2pt}{40}{Autunno.}
\VUsaut{3.0mm}

%%K t04-c1-01-1 p
1. --- I giovani sono cresciuti. Uno di loro, Giovanni, si sposa. Siamo al
\VUgras{pranzo di nozze}, che riunisce intorno alla stessa tavola le famiglie della
sposa e dello sposo.

%%K t04-c1-02-1 p
2. --- Siamo in \VUgras{autunno}, nel mese di \VUgras{settembre}. Il vento freddo di
\VUgras{ottobre} e di \VUgras{novembre} non ha ancora spogliato la campagna delle sue
foglie verdi. L'aria della sera è mite e profumata; per questo si cena sotto la
\VUgras{veranda} \textsuperscript{(8)} che dà sul giardino.

%%K t04-c1-03-1 p
3. --- In lontananza, sulla \VUgras{collina} \textsuperscript{(3)}, si vede la casa dei
genitori di Giovanni, un'elegante \VUgras{villa di campagna} \textsuperscript{(1)}
nascosta nel verde; belle vigne, che danno un vino eccellente, coprono il pendio. Il
vignaiolo se ne prende cura.

%%K t04-c1-04-1 p
4. --- Sopra le viti passa una brigata di \VUgras{pernici} \textsuperscript{(2)},
spaventata dall'arrivo del \VUgras{guardiacaccia} \textsuperscript{(5)}. Lui, con il
\VUgras{fucile} sotto il braccio e il \VUgras{carniere a tracolla}, batte i
\VUgras{cespugli} \textsuperscript{(4)} con il suo \VUgras{cane}
\textsuperscript{(6)}. Sorveglia la tenuta e impedisce ai bracconieri di distruggere
la selvaggina.

%%K t04-c1-05-1 p
5. --- Fuori dalla \VUgras{veranda} sale una \VUgras{vite a pergola}, da cui pendono
pesanti \VUgras{grappoli} \textsuperscript{(7)}.

%%K t04-c1-06-1 p
6. --- I bei fiori d'autunno che ornano la \VUgras{tavola} \textsuperscript{(16)} sono
\VUgras{crisantemi} \textsuperscript{(9)}; il \VUgras{padre} \textsuperscript{(10)}
della sposa se n'è messo uno all'occhiello del frac.

%%K t04-c1-07-1 p
7. --- Il pranzo sta per finire. Il \VUgras{domestico} \textsuperscript{(11)} comincia
a portare i piatti del dolce e fra poco toglierà le varie parti del coperto ---
\VUgras{cucchiai} \textsuperscript{(14)}, \VUgras{forchette} \textsuperscript{(12)},
\VUgras{coltelli} \textsuperscript{(13)}, \VUgras{piatti di portata}
\textsuperscript{(15)} --- che non servono più.

%%K t04-tit-3 sub
\VUpk{10.2pt}{40}{La famiglia.}
\VUsaut{3.0mm}

%%K t04-c1-08-1 p
8. --- Tutti hanno l'aria allegra, e il sorriso con cui la \VUgras{madre}
\textsuperscript{(17)} dello sposo guarda i suoi figli mostra quanto è felice.

%%K t04-c1-09-1 p
9. --- Questo matrimonio unisce due famiglie benestanti della zona. Giovanni, il
\VUgras{marito} \textsuperscript{(18)}, è il \VUgras{figlio} di un grande proprietario
terriero e diventa il \VUgras{genero} di un industriale di successo.

%%K t04-c1-10-1 p
10. --- Gli \VUgras{sposi} sono giovani, eppure sono \VUgras{fidanzati} da molto tempo.
La \VUgras{giovane sposa} \textsuperscript{(19)}, Marta, è incantevole nel suo
\VUgras{abito bianco} ornato di fiori d'arancio.

%%K t04-c1-11-1 p
11. --- Il suo \VUgras{suocero} \textsuperscript{(20)} è felice di avere una
\VUgras{nuora} così cara, e la \VUgras{suocera} \textsuperscript{(21)} di Giovanni
sorride nel vedere la sua \VUgras{figlia} così bella.

%%K t04-c1-12-1 p
12. --- In fondo alla tavola siedono gli altri \VUgras{invitati}
\textsuperscript{(22)}: \VUgras{parenti}, \VUgras{amici}, \VUgras{cugini}. Un
\VUgras{capocameriere} \textsuperscript{(23)}, con il tovagliolo sul braccio, li serve
con prontezza.

%%K t04-tit-4 sub
\VUpk{10.2pt}{40}{Gli amici.}
\VUsaut{3.0mm}

%%K t04-c1-13-1 p
13. --- A destra della tavola siede una \VUgras{signorina} \textsuperscript{(24)},
un'\VUgras{amica} della sposa, poi il \VUgras{fratello} della sposa, che è il suo
\VUgras{testimone} \textsuperscript{(25)}. Chiacchiera con la sua
\VUgras{damigella d'onore} \textsuperscript{(26)}, la sorella dello sposo, che con
questo matrimonio diventa \VUgras{cognata} di Marta. È una giovane donna incantevole.
Ha bei \VUgras{gioielli}, una \VUgras{collana di perle} e un \VUgras{braccialetto}
d'oro.

%%K t04-c1-14-1 p
14. --- Accanto a loro il signore in piedi che alza il bicchiere è un \VUgras{amico}
\textsuperscript{(27)} di famiglia. È uno dei \VUgras{testimoni dello sposo}. Fa un
\VUgras{brindisi}, beve alla \VUgras{salute} degli sposi e augura loro una lunga
felicità. È in abito da cerimonia: frac e \VUgras{scarpe di vernice}
\textsuperscript{(28)}. Accompagna la \VUgras{nonna} \textsuperscript{(29)}, che è
\VUgras{vedova}.

%%K t04-c1-15-1 p
15. --- Il giovane in \VUgras{frac} \textsuperscript{(31)}, con i baffi neri sottili, è
il \VUgras{cognato} \textsuperscript{(30)} di Giovanni. È un ingegnere di prestigio.
È seduto accanto a una \VUgras{signora} \textsuperscript{(33)} molto elegante, la
\VUgras{sorella maggiore} di Marta e madre della bambina graziosa che arriva, al
braccio di suo cugino, per offrire un fiore e \VUgras{augurare} la \VUgras{buonasera}.
Questa signora ha bellissimi \VUgras{orecchini} e un'elegante \VUgras{acconciatura}.

%%K t04-c1-16-1 p
16. --- Il cuginetto, che ha un'aria così buffa con il suo \VUgras{grembiule}
\textsuperscript{(36)} e i suoi \VUgras{calzoni} larghi \textsuperscript{(37)}, fa
pena. È \VUgras{orfano} \textsuperscript{(35)}. Ha perso il padre e la madre da
piccolissimo. Suo zio è il suo \VUgras{tutore} \textsuperscript{(38)} ed è rimasto
scapolo per crescere il povero bambino. È un brav'uomo. È un po' calvo e molto
miope; porta le \VUgras{lenti a molla}.

%%K t04-tit-5 sub
\VUsaut{2.4mm}
\VUpk{10.2pt}{40}{Il dolce.}
\VUsaut{3.0mm}

%%K t04-c1-17-1 p
17. --- Dietro gli invitati, su un tavolo coperto da una \VUgras{tovaglia}, è
apparecchiato il \VUgras{dolce} \textsuperscript{(39)}. In una grande coppa c'è
frutta: \VUgras{pesche} \textsuperscript{(40)}, \VUgras{pere} \textsuperscript{(41)},
\VUgras{mele} \textsuperscript{(42)}, \VUgras{albicocche} \textsuperscript{(43)} e
\VUgras{uva} \textsuperscript{(44)}. Accanto, \VUgras{ananas}
\textsuperscript{(45)}, una bella \VUgras{torta} \textsuperscript{(46)},
\VUgras{bottiglie di liquore} \textsuperscript{(48)} e \VUgras{champagne}
\textsuperscript{(49)}, e anche pile di \VUgras{piatti} \textsuperscript{(47)} puliti.

%%K t04-c1-18-1 p
18. --- Il \VUgras{sommelier} \textsuperscript{(50)} stappa una \VUgras{bottiglia}
\textsuperscript{(52)} di vino vecchio con un \VUgras{cavatappi}
\textsuperscript{(51)}. Il \VUgras{tappo} \textsuperscript{(53)} sembra molto duro da
tirare. Questo sommelier ha un \VUgras{tovagliolo} \textsuperscript{(54)} sul braccio.

%%K t04-c1-19-1 p
19. --- Dopo il pranzo i signori passeranno nella sala da fumo, e le signore andranno a
prepararsi per il ballo.

%%K t04-tit-6 sub
\VUsaut{5.0mm}
\VUcentre{11.4pt}{0}{\textit{Seconda scena.}}
\VUsaut{1.6mm}
\VUpk{11.4pt}{40}{Il compleanno del nonno.}
\VUsaut{2.6mm}

%%K t04-tit-7 sub
\VUpk{10.2pt}{40}{La visita.}
\VUsaut{3.0mm}

%%K t04-c2-01-1 p
1. --- Sono passati dieci anni. I genitori di Giovanni sono invecchiati e vogliono
molto bene ai loro tre nipoti, due \VUgras{nipoti maschi} \textsuperscript{(2)} e una
\VUgras{nipote} \textsuperscript{(3)}. Siccome oggi è il \VUgras{compleanno del nonno},
i bambini sono venuti a fargli gli auguri.

%%K t04-c2-02-1 p
2. --- Il piccolo Pietro è ancora molto piccolo; ha appena tre anni. Vede il
\VUgras{gatto} \textsuperscript{(1)} che gli viene incontro. Vuole accarezzargli il bel
\VUgras{pelo} setoso e la lunga \VUgras{coda}; non gli passa per la testa che sotto
quelle \VUgras{zampe} così graziose ci siano brutte unghie aguzze che potrebbero
graffiarlo facilmente.

%%K t04-c2-03-1 p
3. --- Sua sorella Gabriella, che lo tiene per mano, è molto più seria, e Marcello, con
il suo grande \VUgras{cappotto} \textsuperscript{(4)} e la \VUgras{cintura}
\textsuperscript{(5)} di cuoio, sembra proprio un ometto, malgrado i calzoni corti che
lasciano vedere le sue \VUgras{calze} \textsuperscript{(6)}.

%%K t04-c2-04-1 p
4. --- Il loro padre si è messo il suo pesante \VUgras{cappotto} d'inverno
\textsuperscript{(7)} con il \VUgras{collo di pelliccia} \textsuperscript{(8)}, e in
\VUgras{tasca} \textsuperscript{(9)} ha una sciarpa. Sua moglie ha il cappello di
feltro ornato di \VUgras{piume} \textsuperscript{(10)}, la sua \VUgras{giacca}
\textsuperscript{(11)}, il suo \VUgras{boa} di pelliccia \textsuperscript{(12)} e il
suo \VUgras{manicotto} \textsuperscript{(13)}.

%%K t04-c2-05-1 p
5. --- Fa molto freddo, ora che è arrivato l'inverno. La \VUgras{neve}
\textsuperscript{(19)} è caduta tutto il giorno e i tetti delle case sono tutti
bianchi. Con la \VUgras{notte} il freddo diventa ancora più pungente. In cielo la
\VUgras{luna} \textsuperscript{(17)} e le \VUgras{stelle} \textsuperscript{(18)}
brillano e attraversano il \VUgras{buio}.

%%K t04-tit-8 sub
\VUsaut{4.2mm}
\VUpk{10.2pt}{40}{La malattia.}
\VUsaut{3.0mm}

%%K t04-c2-06-1 p
6. --- Il povero \VUgras{cieco} \textsuperscript{(20)}, che attraversa la piazza davanti
alla \VUgras{farmacia} \textsuperscript{(16)} guidato dal suo \VUgras{barboncino}
\textsuperscript{(21)}, fa molta pena. Giustina, la \VUgras{domestica}
\textsuperscript{(14)}, porta dalla cucina una \VUgras{scodella}
\textsuperscript{(15)} di \VUgras{tisana} bollente, molto zuccherata.

%%K t04-c2-07-1 p
7. --- La \VUgras{nonna} \textsuperscript{(23)}, che vuole prendere dal tavolo i suoi
\VUgras{occhiali} \textsuperscript{(26)} per leggere la \VUgras{ricetta}
\textsuperscript{(27)} che il \VUgras{medico} \textsuperscript{(25)} ha appena
scritto, si alza dalla poltrona appoggiandosi al suo \VUgras{bastone}
\textsuperscript{(24)}.

%%K t04-c2-08-1 p
8. --- Soffre spesso di piccoli \VUgras{disturbi}, \VUgras{giramenti di testa},
\VUgras{raffreddori} e \VUgras{mal di denti}; ha le gambe deboli e gli occhi non più
molto buoni. Con tutto questo le piace prendere in giro il suo vecchio
\VUgras{medico}, che vuole sempre prescriverle una \VUgras{pozione}
\textsuperscript{(28)}. Oggi è molto contenta di avere intorno la sua giovane
famiglia.

%%K t04-c2-09-1 p
9. --- Nella stanza ben chiusa si sta al caldo. Nel \VUgras{camino}
\textsuperscript{(29)} di marmo arde uno \VUgras{splendido fuoco}
\textsuperscript{(30)}, e si sta bene nella \VUgras{luce} morbida della
\VUgras{lampada} \textsuperscript{(31)} appena accesa.

%%K t04-tit-9 sub
\VUsaut{4.2mm}
\VUpk{10.2pt}{40}{Il nonno.}
\VUsaut{3.0mm}

%%K t04-c2-10-1 p
10. --- Perfino il \VUgras{nonno} \textsuperscript{(33)} ha dimenticato di essere stato
malato, che i suoi \VUgras{reumatismi} lo tengono inchiodato alla \VUgras{poltrona}
\textsuperscript{(34)} e che appena ieri ha avuto \VUgras{febbre e svenimenti}. Non
si ricorda più che l'inverno scorso ha avuto una \VUgras{polmonite} e che una
\VUgras{tosse} dura e dolorosa lo ha fatto molto soffrire.

%%K t04-c2-10-2 p
Eppure è guarito a fatica. Adesso sta un po' meglio. Ma anche la gamba, che tiene su
un \VUgras{cuscino} \textsuperscript{(38)} posato su uno \VUgras{sgabello}
\textsuperscript{(39)}, gli fa ancora male.

%%K t04-c2-11-1 p
11. --- Quando è avvolto con cura nella sua calda \VUgras{vestaglia}
\textsuperscript{(35)} con i grandi \VUgras{bottoni} di madreperla
\textsuperscript{(36)}, e quando ha accanto, a leggergli il giornale, la piccola
\VUgras{orfana} \textsuperscript{(40)} con la sua \VUgras{cuffia} bianca, il suo
\VUgras{vestito} azzurro \textsuperscript{(42)} e la sua \VUgras{mantellina}
\textsuperscript{(41)}, allora non si lamenta.

%%K t04-c2-12-1 p
12. --- Ogni anno, quando il \VUgras{calendario} \textsuperscript{(43)} riporta la
\VUgras{data} del primo \VUgras{dicembre}, aspetta con impazienza i suoi
\VUgras{nipoti}, perché sa che non dimenticheranno quella \VUgras{data} importante.

%%K t04-c2-13-1 p
13. --- Ricorda con affetto i giorni lontani in cui andava lui stesso a fare gli auguri
al glorioso \VUgras{maresciallo} imperiale, il cui \VUgras{ritratto}
\textsuperscript{(44)} è appeso al muro e che è il \VUgras{bisnonno} dei bambini.
\VUsaut{6.0mm}
\VUfilet{22mm}
